\section{Related Work}

\noindent\textbf{Representation guidance for diffusion transformers.} Recent work shows that diffusion and flow-based generative models learn intermediate representations whose quality is closely tied to generative performance~\citep{mukhopadhyay2024, kim2025, stracke2025, yu2024repa}. REPA improves DiT and SiT training by aligning hidden states with external visual representations from pretrained encoders~\citep{yu2024repa}, and follow-up methods further refine this alignment paradigm through stronger training recipes or more spatially faithful targets~\citep{wang2025haste,leng2025repae,singh2025irepa}. These methods are effective, but their dependence on external teachers introduces additional computation, architectural coupling, and limited portability to domains where strong pretrained encoders are unavailable~\citep{oquab2023dinov2,yu2024repa}.

\noindent\textbf{Self-contained internal alignment.} To remove external encoders, SRA aligns earlier or noisier features with later or cleaner features from an EMA teacher built from the diffusion model itself~\citep{jiang2025sra}. LayerSync similarly uses deeper, more informative internal layers as intrinsic guidance for weaker shallow layers~\citep{haghighi2026layersync}. Complementary work also suggests that excessive layer similarity can hurt specialization, motivating regularizers that preserve diversity across blocks~\citep{wang2025dispersive,yang2026diversedit}. ReSync follows the self-contained direction of SRA and LayerSync, but replaces direct pairwise alignment with a learned transition model over depth. Rather than only encouraging two layers to be similar, ReSync learns how representations should move from $h_a$ to $h_b$ and enforces that these transitions compose consistently across layers.
